% SPDX-License-Identifier: CC-BY-NC-ND-4.0
% !TEX root = main.tex

\section{随伴}
\begin{definition} \label{def: compositeFunctor}
$\mathscr{C},\ \mathscr{D},\ \mathscr{E}$を圏とする.\ $F \colon \mathscr{C} \to \mathscr{D}$と$G \colon \mathscr{D} \to \mathscr{E}$を共変関手とする.\ このとき\textbf{合成関手}(composite functor)$GF \colon \mathscr{C} \to \mathscr{D}$が次のようにして定まる.
\begin{itemize}
	\item $A \in \Ob(\mathscr{C})$に対して$GF(A) \coloneq G(F(A))$とする.
	\item $f \in \Mor(\mathscr{C})$に対して$GF(f) \coloneq G(F(f))$とする.
\end{itemize}
\end{definition}

\begin{definition} \label{def: leftWhiskering}
$\mathscr{C},\ \mathscr{D},\ \mathscr{E}$を圏,\ $F,\ G \in \mathscr{D}^{\mathscr{C}}$,\ $H \in \mathscr{E}^{\mathscr{D}}$,\ $\eta \colon F \rightarrow G$とする.\ 自然変換$H\eta \colon HF \Rightarrow HG$を$X \in \mathscr{C}$に対して$(H\eta)_X \coloneq H(\eta_X)$として定める.\ $H \eta$を$H$による$\eta$への\textbf{左ヒゲ結合}(left whiskering)という.
\end{definition}

\begin{definition} \label{def: rightWhiskering}
$\mathscr{C},\ \mathscr{D},\ \mathscr{E}$を圏,\ $F,\ G \in \mathscr{E}^{\mathscr{D}}$,\ $H \in \mathscr{D}^{\mathscr{C}}$,\ $\eta \colon F \rightarrow G$とする.\ 自然変換$\eta	H \colon FH \Rightarrow GH$を$X \in \mathscr{C}$に対して$(\eta H)_X \coloneq \eta_{H(X)}$として定める.\ $\eta H$を$H$による$\eta$への\textbf{右ヒゲ結合}(right whiskering)という.
\end{definition}

\begin{definition} \label{def: Adjunction}
$\mathscr{C},\ \mathscr{D}$を圏とする.\ $F \colon \mathscr{C} \to \mathscr{D}$と$G \colon \mathscr{D} \to \mathscr{C}$を共変関手とする.\ 自然変換$\eta \colon \id_{\mathscr{C}} \Rightarrow GF$と自然変換$\varepsilon \colon FG \Rightarrow \id_{\mathscr{D}}$で$\varepsilon F \circ F \eta = \id_{F}$かつ$G \varepsilon \circ \eta G = \id_{G}$を満たすものが存在するとき$F \dashv G \coloneq (\eta, \varepsilon)$としこれを\textbf{随伴}(adjunction)という.\ 随伴$F \dashv G$が存在するとき$F$は$G$の\textbf{左随伴}(left adjoint),\ $G$は$F$の\textbf{右随伴}(right adjoint)であるといい$\eta$を\textbf{単位}(unit),\ $\varepsilon$を\textbf{余単位}(counit)という.

% https://q.uiver.app/#q=WzAsOCxbMCwwLCJcXG1hdGhybXtkb219KGYpIl0sWzAsMSwiXFxtYXRocm17Y29kfShmKSJdLFsxLDAsIkcoRihcXG1hdGhybXtkb219KGYpKSkiXSxbMSwxLCJHKEYoXFxtYXRocm17Y29kfShmKSkpIl0sWzMsMCwiRihHKFxcbWF0aHJte2RvbX0oZykpKSJdLFszLDEsIkYoRyhcXG1hdGhybXtjb2R9KGcpKSkiXSxbNCwwLCJcXG1hdGhybXtkb219KGcpIl0sWzQsMSwiXFxtYXRocm17Y29kfShnKSJdLFswLDIsIlxcZXRhX3tcXG1hdGhybXtkb219KGYpfSJdLFsxLDMsIlxcZXRhX3tcXG1hdGhybXtjb2R9KGYpfSJdLFsyLDMsIkcoRihmKSkiXSxbMCwxLCJmIiwyXSxbNCw2LCJcXHZhcmVwc2lsb25fe1xcbWF0aHJte2RvbX0oZyl9Il0sWzUsNywiXFx2YXJlcHNpbG9uX3tcXG1hdGhybXtjb2R9KGcpfSJdLFs0LDUsIkYoRyhnKSkiLDJdLFs2LDcsImciXV0=
\[\begin{tikzcd}
	{\dom(f)} & {G(F(\dom(f)))} && {F(G(\dom(g)))} & {\dom(g)} \\
	{\cod(f)} & {G(F(\cod(f)))} && {F(G(\cod(g)))} & {\cod(g)}
	\arrow["{\eta_{\dom(f)}}", from=1-1, to=1-2]
	\arrow["f"', from=1-1, to=2-1]
	\arrow["{G(F(f))}", from=1-2, to=2-2]
	\arrow["{\varepsilon_{\dom(g)}}", from=1-4, to=1-5]
	\arrow["{F(G(g))}"', from=1-4, to=2-4]
	\arrow["g", from=1-5, to=2-5]
	\arrow["{\eta_{\cod(f)}}", from=2-1, to=2-2]
	\arrow["{\varepsilon_{\cod(g)}}", from=2-4, to=2-5]
\end{tikzcd}\]

% https://q.uiver.app/#q=WzAsNixbMCwwLCJGKEEpIl0sWzEsMCwiRihHKEYoQSkpKSJdLFsyLDAsIkYoQSkiXSxbMCwxLCJHKEIpIl0sWzEsMSwiRyhGKEcoQikpKSJdLFsyLDEsIkcoQikiXSxbMCwxLCJGKFxcZXRhX0EpIiwyXSxbMSwyLCJcXHZhcmVwc2lsb25fe0YoQSl9IiwyXSxbMCwyLCJcXG1hdGhybXtpZH1fe0YoQSl9IiwwLHsiY3VydmUiOi0zfV0sWzMsNCwiXFxldGFfe0coQil9Il0sWzQsNSwiRyhcXHZhcmVwc2lsb25fe0J9KSJdLFszLDUsIlxcbWF0aHJte2lkfV97RyhCKX0iLDIseyJjdXJ2ZSI6M31dXQ==
\[\begin{tikzcd}
	{F(A)} & {F(G(F(A)))} & {F(A)} \\
	{G(B)} & {G(F(G(B)))} & {G(B)}
	\arrow["{F(\eta_A)}"', from=1-1, to=1-2]
	\arrow["{\id_{F(A)}}", curve={height=-18pt}, from=1-1, to=1-3]
	\arrow["{\varepsilon_{F(A)}}"', from=1-2, to=1-3]
	\arrow["{\eta_{G(B)}}", from=2-1, to=2-2]
	\arrow["{\id_{G(B)}}"', curve={height=18pt}, from=2-1, to=2-3]
	\arrow["{G(\varepsilon_{B})}", from=2-2, to=2-3]
\end{tikzcd}\]
\end{definition}

\begin{definition} \label{def: ProductCategory}
$\mathscr{C}, \mathscr{D}$を圏とする.\ \textbf{積圏}(product category)$\mathscr{C} \times \mathscr{D}$は$A \in \Ob(\mathscr{C}), B \in \Ob(\mathscr{D})$の組$(A, B)$を対象とし$f \in \Mor(\mathscr{C}), g \in \Mor(\mathscr{D})$の組$(f, g)$を射とする.\ 合成は$(f_2, g_2) \circ (f_1, g_1) \coloneq (f_2 \circ f_1, g_2 \circ g_1)$で定める.\ 恒等射は$\id_{(A,B)} \coloneq (\id_A, \id_B)$で定める.
\end{definition}

\begin{definition} \label{def: Bifunctor}
$\mathscr{C}, \mathscr{D}, \mathscr{E}$を圏とする.\ $\mathscr{E}^{\mathscr{C} \times \mathscr{D}}$の対象を\textbf{双関手}(bifunctor)と呼ぶ.
\end{definition}

\begin{proposition} \label{prop: adjunction_iff_hom_natural_iso}
$\mathscr{C}, \mathscr{D}$を局所小圏とし$F \colon \mathscr{C} \to \mathscr{D}, G \colon \mathscr{D} \to \mathscr{C}$を関手とする.\ $\Mor_\mathscr{D}(F(-), -) \colon \mathscr{C}^{\op} \times \mathscr{D} \to \mathbf{Set}, \Mor_\mathscr{C}(-, G(-)) \colon \mathscr{C}^{\op} \times \mathscr{D} \to \mathbf{Set}$を双関手とする.\ このとき次が成り立つ.
\begin{center}
	自然同型$\Phi \colon \Mor_\mathscr{D}(F(-), -) \Rightarrow \Mor_\mathscr{C}(-, G(-))$が存在$\iff$$F \dashv G$が存在
\end{center}
\end{proposition}

\begin{proof}
自然同型$\Phi \colon \Mor_\mathscr{D}(F(-), -) \Rightarrow \Mor_\mathscr{C}(-, G(-))$が存在するとする.\ 任意の$X \in \mathscr{C},\ Y \in \mathscr{D}$に対して$\eta_X \coloneq \Phi_{(X, F(X))}(\id_{F(X)}),\ \varepsilon_Y \coloneq \Phi^{-1}_{(G(Y), Y)}(\id_{G(Y)})$とする.

$X,X' \in \Ob(\mathscr{C}),\ Y,Y' \in \Ob(\mathscr{D}),\ u \in \Mor_{\mathscr{C}}(X', X),\ v \in \Mor_{\mathscr{D}}(Y, Y'),\ f \in \Mor_{\mathscr{D}}(F(X), Y),\ g \in \Mor_{\mathscr{C}}(X, G(Y))$とする.\ このとき$\Mor_{\mathscr{D}}(F(u), v)(f) = v \circ f \circ F(u)$かつ$\Mor_{\mathscr{C}}(u, G(v))(g) \coloneq G(v) \circ g \circ u$である.\ $\Phi$は自然同型なので次の図式が可換になる.
% https://q.uiver.app/#q=WzAsNCxbMCwwLCJcXG1hdGhybXtNb3J9X3tcXG1hdGhzY3J7RH19KEYoWCksWSkiXSxbMSwwLCJcXG1hdGhybXtNb3J9X3tcXG1hdGhzY3J7Q319KFgsRyhZKSkiXSxbMCwxLCJcXG1hdGhybXtNb3J9X3tcXG1hdGhzY3J7RH19KEYoWCcpLFknKSJdLFsxLDEsIlxcbWF0aHJte01vcn1fe1xcbWF0aHNjcntDfX0oWCcsRyhZJykpIl0sWzAsMiwidiBcXGNpcmMgKC0pIFxcY2lyYyBGKHUpIiwyXSxbMCwxLCJcXFBoaV97KFgsIFkpfSIsMCx7Im9mZnNldCI6LTF9XSxbMSwzLCJHKHYpIFxcY2lyYyAoLSkgXFxjaXJjIHUiXSxbMiwzLCJcXFBoaV97KFgnLCBZJyl9IiwwLHsib2Zmc2V0IjotMX1dLFsxLDAsIlxcUGhpXnstMX1feyhYLCBZKX0iLDAseyJvZmZzZXQiOi0xfV0sWzMsMiwiXFxQaGleey0xfV97KFgnLCBZJyl9IiwwLHsib2Zmc2V0IjotMX1dXQ==
\[\begin{tikzcd}
	{\mathrm{Mor}_{\mathscr{D}}(F(X),Y)} & {\mathrm{Mor}_{\mathscr{C}}(X,G(Y))} \\
	{\mathrm{Mor}_{\mathscr{D}}(F(X'),Y')} & {\mathrm{Mor}_{\mathscr{C}}(X',G(Y'))}
	\arrow["{\Phi_{(X, Y)}}", shift left, from=1-1, to=1-2]
	\arrow["{v \circ (-) \circ F(u)}"', from=1-1, to=2-1]
	\arrow["{\Phi^{-1}_{(X, Y)}}", shift left, from=1-2, to=1-1]
	\arrow["{G(v) \circ (-) \circ u}", from=1-2, to=2-2]
	\arrow["{\Phi_{(X', Y')}}", shift left, from=2-1, to=2-2]
	\arrow["{\Phi^{-1}_{(X', Y')}}", shift left, from=2-2, to=2-1]
\end{tikzcd}\]

したがって上の可換図式により任意の$X, X' \in \Ob(\mathscr{C})$と$f \in \Mor_{\mathscr{C}}(X, X')$に対して次の図式が可換になるので$\Phi_{(X,F(X'))}(F(f)) = \Phi_{(X,F(X'))}(\id_{F(X')} \circ \id_F(X') \circ F(f)) = G(\id_{F(X')}) \circ \Phi_{(X,F(X'))}(\id_{F(X')}) \circ f = \Phi_{(X,F(X'))}(\id_{F(X')}) \circ f = \eta_{X'} \circ f$が成り立つ.

% https://q.uiver.app/#q=WzAsNCxbMCwwLCJcXG1hdGhybXtNb3J9X3tcXG1hdGhzY3J7RH19KEYoWCcpLEYoWCcpKSJdLFsxLDAsIlxcbWF0aHJte01vcn1fe1xcbWF0aHNjcntDfX0oWCcsR0YoWCcpKSJdLFswLDEsIlxcbWF0aHJte01vcn1fe1xcbWF0aHNjcntEfX0oRihYKSxGKFgnKSkiXSxbMSwxLCJcXG1hdGhybXtNb3J9X3tcXG1hdGhzY3J7Q319KFgsR0YoWCcpKSJdLFswLDIsIlxcbWF0aHJte2lkfV97RihYJyl9IFxcY2lyYyAoLSkgXFxjaXJjIEYoZikiLDJdLFswLDEsIlxcUGhpX3soWCcsIEYoWCcpKX0iXSxbMSwzLCJHKFxcbWF0aHJte2lkfV97RihYJyl9KSBcXGNpcmMgKC0pIFxcY2lyYyBmIl0sWzIsMywiXFxQaGlfeyhYLCBGKFgnKSl9Il1d
\[\begin{tikzcd}
	{\mathrm{Mor}_{\mathscr{D}}(F(X'),F(X'))} & {\mathrm{Mor}_{\mathscr{C}}(X',GF(X'))} \\
	{\mathrm{Mor}_{\mathscr{D}}(F(X),F(X'))} & {\mathrm{Mor}_{\mathscr{C}}(X,GF(X'))}
	\arrow["{\Phi_{(X', F(X'))}}", from=1-1, to=1-2]
	\arrow["{\mathrm{id}_{F(X')} \circ (-) \circ F(f)}"', from=1-1, to=2-1]
	\arrow["{G(\mathrm{id}_{F(X')}) \circ (-) \circ f}", from=1-2, to=2-2]
	\arrow["{\Phi_{(X, F(X'))}}", from=2-1, to=2-2]
\end{tikzcd}\]

また次の図式も任意の$X, X' \in \Ob(\mathscr{C})$と$f \in \Mor_{\mathscr{C}}(X, X')$に対して可換になるので$\Phi_{(X,F(X'))}(F(f)) = \Phi_{(X,F(X'))}(F(f) \circ \id_{F(X)} \circ F(\id_X)) = G(F(f)) \circ \Phi_{(X,F(X))}(\id_{F(X)}) \circ \id_X = G(F(f)) \circ \Phi_{(X,F(X))}(\id_{F(X)}) = G(F(f)) \circ \eta_X$が成り立つ.
% https://q.uiver.app/#q=WzAsNCxbMCwwLCJcXG1hdGhybXtNb3J9X3tcXG1hdGhzY3J7RH19KEYoWCksRihYKSkiXSxbMSwwLCJcXG1hdGhybXtNb3J9X3tcXG1hdGhzY3J7Q319KFgsR0YoWCkpIl0sWzAsMSwiXFxtYXRocm17TW9yfV97XFxtYXRoc2Nye0R9fShGKFgpLEYoWCcpKSJdLFsxLDEsIlxcbWF0aHJte01vcn1fe1xcbWF0aHNjcntDfX0oWCcsR0YoWCcpKSJdLFswLDIsIkYoZikgXFxjaXJjICgtKSBcXGNpcmMgRihcXG1hdGhybXtpZH1fWCkiLDJdLFswLDEsIlxcUGhpX3soWCwgRihYKSl9Il0sWzEsMywiRyhGKGYpKSBcXGNpcmMgKC0pIFxcY2lyYyBcXG1hdGhybXtpZH1fWCJdLFsyLDMsIlxcUGhpX3soWCwgRihYJykpfSJdXQ==
\[\begin{tikzcd}
	{\mathrm{Mor}_{\mathscr{D}}(F(X),F(X))} & {\mathrm{Mor}_{\mathscr{C}}(X,GF(X))} \\
	{\mathrm{Mor}_{\mathscr{D}}(F(X),F(X'))} & {\mathrm{Mor}_{\mathscr{C}}(X',GF(X'))}
	\arrow["{\Phi_{(X, F(X))}}", from=1-1, to=1-2]
	\arrow["{F(f) \circ (-) \circ F(\mathrm{id}_X)}"', from=1-1, to=2-1]
	\arrow["{G(F(f)) \circ (-) \circ \mathrm{id}_X}", from=1-2, to=2-2]
	\arrow["{\Phi_{(X, F(X'))}}", from=2-1, to=2-2]
\end{tikzcd}\]

ゆえに$\eta_{X'} \circ f = G(F(f)) \circ \eta_X$が成り立ち次の図式が可換になるので$\eta \coloneq (\eta_X)_{X \in \Ob(\mathscr{C})}$は自然変換である.
% https://q.uiver.app/#q=WzAsNCxbMCwwLCJYIl0sWzEsMCwiR0YoWCkiXSxbMSwxLCJHRihYJykiXSxbMCwxLCJYJyJdLFswLDEsIlxcZXRhX1giXSxbMywyLCJcXGV0YV97WCd9Il0sWzAsMywiZiIsMl0sWzEsMiwiR0YoZikiXV0=
\[\begin{tikzcd}
	X & {GF(X)} \\
	{X'} & {GF(X')}
	\arrow["{\eta_X}", from=1-1, to=1-2]
	\arrow["f"', from=1-1, to=2-1]
	\arrow["{GF(f)}", from=1-2, to=2-2]
	\arrow["{\eta_{X'}}", from=2-1, to=2-2]
\end{tikzcd}\]

$\varepsilon$が自然変換であることも同様にして示せる.

次の図式も任意の$X \in \Ob(\mathscr{C})$に対して可換になるので$\id_{F(X)} = \Phi^{-1}_{(X, F(X))}(\Phi_{(X, F(X))}(\id_{F(X)})) = \Phi^{-1}_{(X, F(X))}(\eta_X) = \Phi^{-1}_{(X, F(X))}(G(\id_{F(X)}) \circ \id_{GF(X)} \circ \eta_X) = \id_{F(X)} \circ \Phi^{-1}_{(GF(X), F(X))}(\id_{GF(X)}) \circ F(\eta_X) = \varepsilon_{F(X)} \circ F(\eta_X)$が成り立つ.\ つまり$\id_F = \varepsilon F \circ F \eta$である.
% https://q.uiver.app/#q=WzAsNCxbMCwwLCJcXG1hdGhybXtNb3J9X3tcXG1hdGhzY3J7RH19KEZHRihYKSxGKFgpKSJdLFsxLDAsIlxcbWF0aHJte01vcn1fe1xcbWF0aHNjcntDfX0oR0YoWCksR0YoWCkpIl0sWzAsMSwiXFxtYXRocm17TW9yfV97XFxtYXRoc2Nye0R9fShGKFgpLEYoWCkpIl0sWzEsMSwiXFxtYXRocm17TW9yfV97XFxtYXRoc2Nye0N9fShYLEdGKFgpKSJdLFswLDIsIlxcbWF0aHJte2lkfV97RihYKX0gXFxjaXJjICgtKSBcXGNpcmMgRihcXGV0YV9YKSIsMl0sWzEsMywiRyhcXG1hdGhybXtpZH1fe0YoWCl9KSBcXGNpcmMgKC0pIFxcY2lyYyBcXGV0YV9YIl0sWzEsMCwiXFxQaGleey0xfV97KEdGKFgpLCBGKFgpKX0iLDJdLFszLDIsIlxcUGhpXnstMX1feyhYLCBGKFgpKX0iLDJdXQ==
\[\begin{tikzcd}
	{\mathrm{Mor}_{\mathscr{D}}(FGF(X),F(X))} & {\mathrm{Mor}_{\mathscr{C}}(GF(X),GF(X))} \\
	{\mathrm{Mor}_{\mathscr{D}}(F(X),F(X))} & {\mathrm{Mor}_{\mathscr{C}}(X,GF(X))}
	\arrow["{\mathrm{id}_{F(X)} \circ (-) \circ F(\eta_X)}"', from=1-1, to=2-1]
	\arrow["{\Phi^{-1}_{(GF(X), F(X))}}"', from=1-2, to=1-1]
	\arrow["{G(\mathrm{id}_{F(X)}) \circ (-) \circ \eta_X}", from=1-2, to=2-2]
	\arrow["{\Phi^{-1}_{(X, F(X))}}"', from=2-2, to=2-1]
\end{tikzcd}\]

同様にして$\id_G = G \varepsilon \circ \eta G$が示せるので$(\eta, \varepsilon)$は随伴である.\ したがって随伴$F \dashv G$が存在する.

$F \dashv G$が存在するとする.\ $(\eta, \varepsilon) \coloneq F \dashv G$とすると$\eta \colon \id_{\mathscr{C}} \Rightarrow GF,\ \varepsilon \colon FG \Rightarrow \id_{\mathscr{D}}$であり$\varepsilon F \circ F \eta = \id_F$かつ$G \varepsilon \circ \eta G = \id_G$が成り立つ.\ $X \in \mathscr{C},\ Y \in \mathscr{D}$に対して$\Phi_{(X,Y)} \colon \Mor_{\mathscr{D}}(F(X), Y) \to \Mor_{\mathscr{C}}(X, G(Y))$を$\Phi_{(X,Y)}(f) \coloneq G(f) \circ \eta_X$で定め$\Psi_{(X,Y)} \colon \Mor_{\mathscr{C}}(X, G(Y)) \to \Mor_{\mathscr{D}}(F(X), Y)$を$\Psi_{(X,Y)}(g) \coloneq \varepsilon_Y \circ F(g)$で定める.

任意の$f \in \Mor_{\mathscr{D}}(F(X), Y)$に対して$\Psi_{(X,Y)}(\Phi_{(X,Y)}(f)) = \Psi_{(X,Y)}(G(f) \circ \eta_X) = \varepsilon_Y \circ F(G(f) \circ \eta_X) = \varepsilon_Y \circ FG(f) \circ (F \eta)_X$となる.\ $\varepsilon$は自然変換なので次の図式が可換になる.
% https://q.uiver.app/#q=WzAsNCxbMCwwLCJGR0YoWCkiXSxbMSwwLCJGKFgpIl0sWzEsMSwiWSJdLFswLDEsIkZHKFkpIl0sWzAsMSwiXFx2YXJlcHNpbG9uX3tGKFgpfSJdLFszLDIsIlxcdmFyZXBzaWxvbl97WX0iXSxbMSwyLCJmIl0sWzAsMywiRkcoZikiLDJdXQ==
\[\begin{tikzcd}
	{FGF(X)} & {F(X)} \\
	{FG(Y)} & Y
	\arrow["{\varepsilon_{F(X)}}", from=1-1, to=1-2]
	\arrow["{FG(f)}"', from=1-1, to=2-1]
	\arrow["f", from=1-2, to=2-2]
	\arrow["{\varepsilon_{Y}}", from=2-1, to=2-2]
\end{tikzcd}\]
したがって$\varepsilon_Y \circ FG(f) \circ (F \eta)_X = f \circ \varepsilon_{F(X)} \circ (F \eta)_X = f \circ (\id_F)_X = f$となる.

任意の$g \in \Mor_{\mathscr{C}}(X, G(Y))$に対して$\Phi_{(X,Y)}(\Psi_{(X,Y)}(g)) = \Phi_{(X,Y)}(\varepsilon_Y \circ F(g)) = G(\varepsilon_Y \circ F(g)) \circ \eta_X = G \varepsilon_Y \circ GF(g) \circ \eta_X$となる.\ $\eta$は自然変換なので次の図式が可換になる.
% https://q.uiver.app/#q=WzAsNCxbMCwwLCJYIl0sWzAsMSwiRyhZKSJdLFsxLDAsIkdGKFgpIl0sWzEsMSwiR0ZHKFkpIl0sWzAsMSwiZyIsMl0sWzAsMiwiXFxldGFfWCJdLFsyLDMsIkdGKGcpIl0sWzEsMywiXFxldGFfe0coWSl9Il1d
\[\begin{tikzcd}
	X & {GF(X)} \\
	{G(Y)} & {GFG(Y)}
	\arrow["{\eta_X}", from=1-1, to=1-2]
	\arrow["g"', from=1-1, to=2-1]
	\arrow["{GF(g)}", from=1-2, to=2-2]
	\arrow["{\eta_{G(Y)}}", from=2-1, to=2-2]
\end{tikzcd}\]
したがって$(G \varepsilon)_Y \circ GF(g) \circ \eta_X = (G \varepsilon)_Y \circ (\eta G)_Y \circ g = (\id_G)_Y \circ g = g$となる.

以上により$\Phi_{(X,Y)}$と$\Psi_{(X,Y)}$は互いに逆写像である.

任意の$X,X' \in \mathscr{C},\ Y, Y' \in \mathscr{D}$と任意の$u \in \Mor_{\mathscr{C}}(X', X),\ v \in \Mor_{\mathscr{D}}(Y, Y'),\ f \in \Mor_{\mathscr{D}}(F(X), Y)$に対して$\Phi_{(X',Y')}(v \circ f \circ F(u)) = G(v \circ f \circ F(u)) \circ \eta_{X'} = G(v) \circ G(f) \circ GF(u) \circ \eta_{X'}$となる.\ $\eta$は自然変換なので次の図式が可換になる.
% https://q.uiver.app/#q=WzAsNCxbMCwwLCJYJyJdLFswLDEsIlgiXSxbMSwwLCJHRihYJykiXSxbMSwxLCJHRihYKSJdLFsyLDMsIkdGKHUpIl0sWzAsMSwidSIsMl0sWzAsMiwiXFxldGFfe1gnfSJdLFsxLDMsIlxcZXRhX1giXV0=
\[\begin{tikzcd}
	{X'} & {GF(X')} \\
	X & {GF(X)}
	\arrow["{\eta_{X'}}", from=1-1, to=1-2]
	\arrow["u"', from=1-1, to=2-1]
	\arrow["{GF(u)}", from=1-2, to=2-2]
	\arrow["{\eta_X}", from=2-1, to=2-2]
\end{tikzcd}\]
よって$G(v) \circ G(f) \circ GF(u) \circ \eta_{X'} = G(v) \circ G(f) \circ \eta_X \circ u = G(v) \circ \Phi_{(X, Y)}(f) \circ u$である.

以上より$\Phi \colon \Mor_{\mathscr{D}}(F(-), -) \Rightarrow \Mor_{\mathscr{C}}(-, G(-))$は自然同型である.

\end{proof}

\begin{proposition} \label{prop: complete_iff_constant_diagram_has_right_adjoint}
局所小圏$\mathscr{C}$が小圏$\mathscr{J}$を形とする極限を持つことと定図式関手$\Delta \colon \mathscr{C}^{\mathscr{D}} \to \mathscr{C}$が右随伴関手を持つことは同値である.
\end{proposition}

\begin{proof}
任意の$F \in \mathscr{C}^{\mathscr{J}}$に対して$ \varprojlim_{\mathscr{J}}F$が存在するとする.\ $ (\varprojlim_{\mathscr{J}}F,\pi_F) \coloneq \varprojlim_{\mathscr{J}}F$とする.\ $A \in \Ob(\mathscr{C})$と$F \in \mathscr{C}^{\mathscr{J}}$に対して$\Mor_{\mathscr{C}^{\mathscr{J}}}(\Delta(A), F)$は定関手$C_A$から$F$への自然変換全体の集合であり写像$\phi \colon \Mor_{\mathscr{C}^{\mathscr{J}}}(\Delta(A), F) \to \Cone(A,F)$を$\phi(\eta) \coloneq (A, \eta)$と定め写像$\psi \colon \Cone(A,F) \to \Mor_{\mathscr{C}^{\mathscr{J}}}(\Delta(A), F)$を$\psi(A, \eta) \coloneq \eta$と定めれば$\phi$と$\psi$は互いに逆写像となり$\Mor_{\mathscr{C}^{\mathscr{J}}}(\Delta(A), F) \cong \Cone(A,F)$である.

$\Phi_A \colon \Mor_{\mathscr{C}}(A, \varprojlim_{\mathscr{J}}F) \to \Cone(A, F)$を$\Phi_A(h) \coloneq (A, \pi \circ \Delta(h))$で定める.\ 任意の$(A, \alpha) \in \Cone(A, F)$に対して$\displaystyle (A, (A, \alpha)) \in \int_{\mathscr{C}^{\op}} \Cone(-, F)$である.\ $(\varprojlim_{\mathscr{J}}F, (\varprojlim_{\mathscr{J}}F, \pi_F))$は$\displaystyle \int_{\mathscr{C}^{\op}} \Cone(-, F)$の終対象なので射$h \colon (A, (A, \alpha)) \to (\varprojlim_{\mathscr{J}}F, (\varprojlim_{\mathscr{J}}F, \pi_F))$がただ1つ存在する.\ $\Delta(h) \colon C_A \Rightarrow C_L$と$\pi_F \colon C_L \Rightarrow F$の合成$\pi_F \circ \Delta(h) \colon C_A \Rightarrow F$は下の図式が可換になることから$\alpha$と等しい.\ つまり$\Phi_A(h) = (A, \alpha)$となる$h$がただ1つ存在するので$\Phi_A$は全単射である.

% https://q.uiver.app/#q=WzAsMTAsWzUsMCwiQ19MKFgpIl0sWzYsMCwiRihYKSJdLFs1LDEsIkNfTChZKSJdLFs2LDEsIkYoWSkiXSxbNCwwLCJDX0EoWCkiXSxbNCwxLCJDX0EoWSkiXSxbMCwwLCJBIl0sWzEsMCwiTCJdLFsyLDAsIkYoWCkiXSxbMiwxLCJGKFkpIl0sWzAsMSwiXFxwaV9YIl0sWzIsMywiXFxwaV9ZIl0sWzAsMiwiQ19MKGYpIiwxXSxbMSwzLCJGKGYpIl0sWzQsMCwiXFxEZWx0YShoKV9YIl0sWzUsMiwiXFxEZWx0YShoKV9ZIl0sWzQsNSwiQ19BKGYpIiwyXSxbNCwxLCJcXGFscGhhX1giLDAseyJjdXJ2ZSI6LTR9XSxbNSwzLCJcXGFscGhhX1kiLDAseyJjdXJ2ZSI6NH1dLFs2LDcsImgiXSxbNyw4LCJcXHBpX1giXSxbOCw5LCJGKGYpIl0sWzcsOSwiXFxwaV9ZIiwyXSxbNiw4LCJcXGFscGhhX1giLDAseyJjdXJ2ZSI6LTN9XSxbNiw5LCJcXGFscGhhX1kiLDIseyJjdXJ2ZSI6Mn1dLFsyMSwxNiwiIiwwLHsic2hvcnRlbiI6eyJzb3VyY2UiOjMwLCJ0YXJnZXQiOjQwfSwic3R5bGUiOnsidGFpbCI6eyJuYW1lIjoiYXJyb3doZWFkIn19fV1d
\[\begin{tikzcd}
	A & L & {F(X)} && {C_A(X)} & {C_L(X)} & {F(X)} \\
	&& {F(Y)} && {C_A(Y)} & {C_L(Y)} & {F(Y)}
	\arrow["h", from=1-1, to=1-2]
	\arrow["{\alpha_X}", curve={height=-18pt}, from=1-1, to=1-3]
	\arrow["{\alpha_Y}"', curve={height=12pt}, from=1-1, to=2-3]
	\arrow["{\pi_X}", from=1-2, to=1-3]
	\arrow["{\pi_Y}"', from=1-2, to=2-3]
	\arrow[""{name=0, anchor=center, inner sep=0}, "{F(f)}", from=1-3, to=2-3]
	\arrow["{\Delta(h)_X}", from=1-5, to=1-6]
	\arrow["{\alpha_X}", curve={height=-24pt}, from=1-5, to=1-7]
	\arrow[""{name=1, anchor=center, inner sep=0}, "{C_A(f)}"', from=1-5, to=2-5]
	\arrow["{\pi_X}", from=1-6, to=1-7]
	\arrow["{C_L(f)}"{description}, from=1-6, to=2-6]
	\arrow["{F(f)}", from=1-7, to=2-7]
	\arrow["{\Delta(h)_Y}", from=2-5, to=2-6]
	\arrow["{\alpha_Y}", curve={height=24pt}, from=2-5, to=2-7]
	\arrow["{\pi_Y}", from=2-6, to=2-7]
	\arrow[between={0.3}{0.6}, Rightarrow, nfold, 2tail reversed, from=0, to=1]
\end{tikzcd}\]

任意の$f \in \Mor_{\mathscr{C}}(A, B)$と$h \in \Mor_{\mathscr{C}}(B, L)$に対して$\Phi_A(h \circ f) = (A, \pi \circ \Delta(h \circ f)) = (A, \pi \circ \Delta(h) \circ \Delta(f)) = \Cone(f, F)(B, \pi \circ \Delta(h)) = \Cone(f, F)(\Phi_B(h))$なので次の図式が可換になる.

% https://q.uiver.app/#q=WzAsNCxbMCwwLCJcXG1hdGhybXtNb3J9X3tcXG1hdGhzY3J7Q319KEIsIFxcdmFycHJvamxpbV97XFxtYXRoc2Nye0p9fUYpIl0sWzEsMCwiXFxtYXRocm17Q29uZX0oQixGKSJdLFsxLDEsIlxcbWF0aHJte0NvbmV9KEEsRikiXSxbMCwxLCJcXG1hdGhybXtNb3J9X3tcXG1hdGhzY3J7Q319KEEsIFxcdmFycHJvamxpbV97XFxtYXRoc2Nye0p9fUYpIl0sWzEsMiwiXFxtYXRocm17Q29uZX0oZixGKSJdLFswLDMsIi0gXFxjaXJjIGYiLDJdLFswLDEsIlxcUGhpX0IiXSxbMywyLCJcXFBoaV9BIl1d
\[\begin{tikzcd}
	{\Mor_{\mathscr{C}}(B, \varprojlim_{\mathscr{J}}F)} & {\Cone(B,F)} \\
	{\Mor_{\mathscr{C}}(A, \varprojlim_{\mathscr{J}}F)} & {\Cone(A,F)}
	\arrow["{\Phi_B}", from=1-1, to=1-2]
	\arrow["{- \circ f}"', from=1-1, to=2-1]
	\arrow["{\Cone(f,F)}", from=1-2, to=2-2]
	\arrow["{\Phi_A}", from=2-1, to=2-2]
\end{tikzcd}\]

したがって$\Phi \colon \Mor_{\mathscr{C}}(-,\varprojlim_{\mathscr{J}}F) \Rightarrow \Cone(-,F)$は自然同型である.\ ゆえに$\Mor_{\mathscr{C}}(-,\varprojlim_{\mathscr{J}}F) \cong \Cone(-,F) \cong \Mor_{\mathscr{C}^{\mathscr{J}}}(\Delta(-), F)$.

自然同型$\Phi \colon \Mor_{\mathscr{C}}(-,\varprojlim_{\mathscr{J}}F) \Rightarrow \Mor_{\mathscr{C}^{\mathscr{J}}}(\Delta(-), F)$が存在すると仮定する.\ 米田の補題により全単射$\phi \colon \Nat(\Mor_{\mathscr{C}}(-, \varprojlim_{\mathscr{J}}F), \Cone(-, F)) \to \Cone(\varprojlim_{\mathscr{J}}F, F)$で$\phi(\alpha) = \alpha_{\varprojlim_{\mathscr{J}}F}(\id_{\varprojlim_{\mathscr{J}}F})$となるものが存在する.\ $\Phi \in \Nat(\Mor_{\mathscr{C}}(-, \varprojlim_{\mathscr{J}}F), \Cone(-, F))$なので$\pi \coloneq \phi(\Phi)$とする.\ 任意の$(A, (A, \alpha)) \in {\displaystyle \int_{\mathscr{C}^{\op}} \Cone(-, F)}$に対して$\Phi_A(h) = \alpha$となる$h \in \Mor_{\mathscr{C}}(A, {\varprojlim_{\mathscr{J}}F})$がただ1つ存在する.\ $\Phi$は自然変換なので$\alpha = \Phi_A(h) = \Cone(h, F)(\varprojlim_{\mathscr{J}}F, \pi) = (\varprojlim_{\mathscr{J}}F, \pi \circ \Delta(h))$が成り立つ.\ よって$h$は$\alpha = \pi \circ \Delta(h)$を満たす$\displaystyle \int_{\mathscr{C}^{\op}} \Cone(-, F)$のただ1つの射$(A, (A, \alpha)) \to (\varprojlim_{\mathscr{J}}F, (\varprojlim_{\mathscr{J}}F, \pi))$なので$(\varprojlim_{\mathscr{J}}F, (\varprojlim_{\mathscr{J}}F, \pi))$は$\displaystyle \int_{\mathscr{C}^{\op}} \Cone(-, F)$の終対象である.

\end{proof}